\documentclass[10pt,a4paper]{amsart}
\usepackage[margin=19mm]{geometry}
\usepackage{amsmath,amssymb,mathtools}
\usepackage[scr=rsfs,cal=euler]{mathalfa}
\usepackage{xfrac}
\usepackage[unicode,hidelinks,bookmarks=false]{hyperref}
\newcommand{\ie}{i.e.\ }
\newcommand{\cf}{cf.\ }
\newcommand{\resp}{respectively\ }
\newcommand{\Subsection}{\subsection}
\providecommand{\nobreakdash}{\nobreak\hbox{-}}
\setlength{\emergencystretch}{2em}
\multlinegap0pt
\usepackage{mathtools}
\usepackage{amssymb}
\usepackage[shortlabels]{enumitem}
\usepackage{xcolor}
\newtheorem{theorem}{Theorem}
\newcommand{\D}{\mathcal{D}}
\newcommand{\Kt}{\widetilde{K}}
\newcommand{\OOh}{\widehat{\mathcal{O}}}
\newcommand{\R}{\mathbb{R}}
\newcommand{\Sigmah}{\widehat{\Sigma}}
\newcommand{\ii}{\mathrm{i}}
\newcommand{\sltwo}{\mathfrak{sl}(2,\R)}
\title{Symmetry Regularization of 1D Generalized Coulomb Problems}
\author{Zhanqiang Bai, Junwei Ma, Guowu Meng}
\date{Source: arXiv:2606.08900v1 (CC BY 4.0)}
\begin{document}
\maketitle

\noindent\emph{Source and licence.} Symmetry Regularization of 1D Generalized Coulomb Problems. Version arXiv:2606.08900v1 (CC BY 4.0), distributed by arXiv under the Creative Commons Attribution 4.0 International licence, http://creativecommons.org/licenses/by/4.0/, as recorded in the arXiv licence metadata for this exact version and shown on the abstract page https://arxiv.org/abs/2606.08900v1. Full licence notice: \texttt{licenses/arxiv-2606.08900-CC-BY-4.0.txt}.\par\smallskip

\noindent\textit{Editorial scope.} Counted unit: the article's theorem thm:iotahat (source0.tex lines 746--777). The article's complete original proof of it is delivered with it (source0.tex lines 779--827, one \texttt{proof} environment of 49 lines). No mathematics is written, repaired or re-derived: the statement and the proof are the article's own lines, in the article's own notation, and no retained line is substituted or re-flowed.  Retained uncounted as the source-local context the proof needs: lines 471--471; lines 538--540; lines 681--684.  Nothing was omitted from any retained span, so the package carries no omission record.  No retained line contains a citation, so the wrapper carries no bibliography and no bibliography entry.  No retained line contains a figure environment, an image-inclusion command or drawing code, so no image is delivered and the package declares no asset dependency.  Editorial formatting only, in addition: an AMS \texttt{amsart} preamble that restates the article's own declarations and macros used by the retained lines, namely the environments \texttt{theorem} on separate counters, together with the standard packages those lines need.  The retained environments print in this document as Theorem 1 in order of appearance, whereas the article prints the same items with the numbering of its own full document.  The complete native source of the article as distributed in the arXiv e-print for this exact version is bundled unchanged as \texttt{sources/202368/source0.tex}.
\subsection{Four concrete models}\label{ssec:four-models}

\begin{equation}\label{eq:hatphin-eigen}
  \hat\pi_{\kappa}(h)\,\hat\phi_{n}\;=\;(\kappa+2n)\,\hat\phi_{n},\qquad n=0,1,2,\ldots.
\end{equation}

\begin{equation}\label{eq:psi-E-Coulomb}
  \psi_{E}(q)\;=\;\mathcal{C}_{\kappa}(k)\,q^{\kappa/2}\,e^{-\ii kq}\;
    {}_{1}F_{1}\!\bigl(\tfrac{\kappa}{2}+\tfrac{\ii}{k};\,\kappa;\,2\ii kq\bigr),
\end{equation}

\begin{theorem}[Quantum regularization]\label{thm:iotahat}
Realize both targets $\OOh_{\pm}$ on the common Whittaker half-line Hilbert space
$L^{2}(\R_{>0},t^{-1}\,\mathrm{d}t)$, equipped respectively with the representations
$\hat\pi_{\kappa}$ for $\OOh_{-}=\D^{+}_{\kappa}$ and $\hat\pi_{(\kappa+1)/2}$
for $\OOh_{+}=\D^{+}_{(\kappa+1)/2}$. Let $E,F$ denote the standard
$\sltwo$-basis elements of Section~\ref{ssec:four-models}.

\smallskip\noindent\textbf{(i) Negative-energy case.}
Set $h=\ii(F-E)$. The map $\psi_{n}\mapsto\hat\phi_{n}$, where
$\{\psi_{n}\}$ is the bound-state eigenbasis of $H_{\kappa}$ on $\Sigmah_{-}$
and $\{\hat\phi_{n}\}$ is the orthonormal $\Kt$-eigenbasis
of~\eqref{eq:hatphin-eigen}, extends by linearity and continuity to a unitary
isomorphism
\[
  \hat\iota_{-}\colon\Sigmah_{-}\xrightarrow{\;\sim\;}\OOh_{-}
\]
satisfying $\hat\iota_{-}^{-1}\,\hat h_{-}\,\hat\iota_{-}=H_{\kappa}$, where
$\hat h_{-}:=-1/\bigl(2(\hat\pi_{\kappa}(h/2))^{2}\bigr)$.

\smallskip\noindent\textbf{(ii) Positive-energy case.}
Let $\psi_{\lambda}:=\lambda^{-1}\,\psi_{E}\big|_{E=1/(2\lambda^{2})}$ denote the
scattering generalized eigenfunction~\eqref{eq:psi-E-Coulomb} reparametrized by
$\lambda=(2E)^{-1/2}>0$, so that
$\int_{0}^{\infty}\overline{\psi_{\lambda}(q)}\,\psi_{\lambda'}(q)\,\mathrm{d}q
=\delta(\lambda-\lambda')$. The map $\psi_{\lambda}\mapsto\sqrt{t}\,\delta(t-\lambda)$
extends to a unitary isomorphism
\[
  \hat\iota_{+}\colon\Sigmah_{+}\xrightarrow{\;\sim\;}\OOh_{+}
\]
satisfying $\hat\iota_{+}^{-1}\,\hat h_{+}\,\hat\iota_{+}=H_{\kappa}$, where
$\hat h_{+}:=1/\bigl(2(-\ii\hat\pi_{(\kappa+1)/2}(E))^{2}\bigr)$.
\end{theorem}

\begin{proof}
The argument in both cases rests on the same principle: two self-adjoint
operators on separable Hilbert spaces that have the same spectral type
(i.e.\ the same spectrum with the same multiplicity) are unitarily
equivalent, and matching (generalized) eigenfunctions with equal
eigenvalues defines the intertwining unitary isomorphism.

\smallskip\noindent\textbf{(i)}
On $\Sigmah_{-}$, the operator
$\hat N_{-}=(-2H_{\kappa})^{-1/2}$ is positive and self-adjoint with
purely discrete, simple spectrum
$\{n+\kappa/2\}_{n=0}^{\infty}$ and orthonormal eigenbasis
$\{\psi_{n}\}$.
On $\OOh_{-}=\D^{+}_{\kappa}$ (Whittaker model), the elliptic Cartan
generator $\hat\pi_{\kappa}(h/2)$ is self-adjoint with the same purely
discrete, simple spectrum $\{n+\kappa/2\}_{n=0}^{\infty}$ and
orthonormal eigenbasis $\{\hat\phi_{n}\}$
(cf.~\eqref{eq:hatphin-eigen}).
Since both eigenbases are complete orthonormal systems indexed by the
same eigenvalues, the map $\psi_{n}\mapsto\hat\phi_{n}$ extends by
linearity and continuity to a unitary isomorphism
$\hat\iota_{-}\colon\Sigmah_{-}\xrightarrow{\sim}\OOh_{-}$
satisfying $\hat\iota_{-}^{-1}\,\hat\pi_{\kappa}(h/2)\,\hat\iota_{-}
=\hat N_{-}$.
Squaring and taking reciprocals gives
$\hat\iota_{-}^{-1}\,\hat h_{-}\,\hat\iota_{-}
=-1/(2\hat N_{-}^{2})=H_{\kappa}$, where
$\hat h_{-}=-1/\bigl(2(\hat\pi_{\kappa}(h/2))^{2}\bigr)$.

\smallskip\noindent\textbf{(ii)}
On $\Sigmah_{+}$, the operator
$\hat N_{+}=(2H_{\kappa})^{-1/2}$ is positive and self-adjoint with
purely continuous, simple spectrum $(0,\infty)$ and Dirac-normalized
generalized eigenfunctions $\{\psi_{\lambda}\}_{\lambda>0}$.
On $\OOh_{+}=\D^{+}_{(\kappa+1)/2}$ (Whittaker model), the self-adjoint
companion $-\ii\hat\pi_{(\kappa+1)/2}(E)$ acts as
multiplication by $t>0$, hence also has purely continuous, simple
spectrum $(0,\infty)$ with Dirac-normalized generalized eigenfunctions
$\{\sqrt{t}\,\delta(t-\lambda)\}_{\lambda>0}$.
Since both operators have the same spectral type, matching generalized
eigenfunctions at each $\lambda$ defines the unitary isomorphism
$\hat\iota_{+}\colon\Sigmah_{+}\xrightarrow{\sim}\OOh_{+}$
satisfying $\hat\iota_{+}^{-1}\,(-\ii\hat\pi_{(\kappa+1)/2}(E))\,\hat\iota_{+}
=\hat N_{+}$.
Squaring and taking reciprocals gives
$\hat\iota_{+}^{-1}\,\hat h_{+}\,\hat\iota_{+}
=1/(2\hat N_{+}^{2})=H_{\kappa}$, where
$\hat h_{+}=1/\bigl(2(-\ii\hat\pi_{(\kappa+1)/2}(E))^{2}\bigr)$.
\end{proof}

\end{document}
